\chapter{Troubleshooting}

\ChapterMeta{
This chapter provides an accelerated debugging playbook: it emphasizes validating artifacts first, followed by investigating common integration failure modes.
}{
Adhering to this playbook drastically shortens iteration times by identifying schema violations, pathing errors, and class-mapping discrepancies before you attempt to interpret potentially flawed metrics.
}{
It is most effective when validators are executed on the exact artifacts slated for evaluation. Pay meticulous attention to the \cmd{--pred-root} parameter and ensure consistent class ordering between the exporter and the evaluator.
}
\ChapterAudience{Read this chapter when results look wrong and you want the shortest route to the first likely failure point.}


\section{Validate early and often}
If evaluation results appear anomalous, your immediate first step must be artifact validation.

To validate the standard predictions schema:
\begin{lstlisting}[language=bash]
python3 tools/validate_predictions.py /path/to/predictions.json --strict
\end{lstlisting}

To validate instance segmentation predictions:
\begin{lstlisting}[language=bash]
python3 tools/validate_instance_segmentation_predictions.py /path/to/preds.json
\end{lstlisting}

\section{Common failure modes}
\begin{itemize}
  \item \textbf{Pathing issues}: Relative paths embedded within the JSON artifact fail to resolve correctly against the specified \cmd{--pred-root}.
  \item \textbf{Class mapping mismatches}: The ordering defined in \path{classes.txt} diverges from the internal mapping utilized by the model or exporter.
  \item \textbf{Backend discrepancies}: Pre-processing or post-processing logic varies subtly across different inference backends.
  \item \textbf{TensorRT environment constraints}: TensorRT execution is strictly pinned to Linux/NVIDIA environments; use the documented container workflow to keep CUDA/TRT/ORT compatibility pinned.
  \item \textbf{CI scope optimization}: Required CI uses a change-scope fast path; docs/metadata-only commits skip heavy test jobs by design.
  \item \textbf{Container workflow expectations}: \path{.github/workflows/container.yml} now builds container-related pull requests for early validation, while image publication remains gated to tag/manual runs. Manual republishes should pass \texttt{release\_tag=vX.Y.Z}; when \texttt{NGC\_API\_KEY} is configured, the same workflow mirrors images to \path{nvcr.io/yolozu/...}.
  \item \textbf{Container trigger scope}: Container checks on pull requests and \texttt{main} pushes stay restricted to container-related paths (workflow/deploy/packaging inputs) to avoid unnecessary runs.
  \item \textbf{Security scan backlog vs real regressions}: \path{.github/workflows/codeql.yml} and \path{.github/workflows/scorecard.yml} report repository posture, not just runtime test failures. Scorecard also runs on pull requests targeting \texttt{main}, so posture regressions should surface before merge rather than only after landing on the default branch.
  \item \textbf{Fuzzing coverage expectations}: \path{.github/workflows/cflite\_pr.yml} and \path{.github/workflows/cflite\_batch.yml} exercise predictions normalization via ClusterFuzzLite. When predictions parsing/canonicalization changes, update \path{.clusterfuzzlite/} and \path{fuzz/} in the same pull request.
  \item \textbf{Self-hosted GPU workflow behavior}: \path{.github/workflows/ngc_test.yml} probes for an idle \texttt{self-hosted + gpu} runner first; if none is online/idle, the workflow exits via a deterministic skip job. If probe API access is denied (\texttt{403}), configure repository secret \texttt{RUNNER\_DISCOVERY\_TOKEN} to enable runner discovery.
  \item \textbf{ONNXRuntime model-load failures}: If ORT reports unsupported ONNX IR versions, regenerate artifacts with a compatible IR version (CI smoke pins IR v11 for compatibility).
\end{itemize}

\section{Further reading and resources}
\begin{itemize}
  \item Tools index: \path{docs/tools_index.md}
  \item External inference documentation: \path{docs/external_inference.md}
  \item Predictions schema specification: \path{docs/predictions_schema.md}
  \item Training and export overview: \path{docs/training_inference_export.md}
  \item CI incident logbook: \path{docs/ci_incidents.md}
  \item Release reliability checklist: \path{docs/release_reliability_checklist.md}
  \item Security and cryptography scope note: \path{docs/security_crypto_scope.md}
\end{itemize}
